Granular-ball-based modeling has been studied for label distribution feature
selection. Qian et al.~\cite{Qian2023GBLDL} proposed a granular-ball
computing-based fuzzy rough set for this task. More recently,
Deng et al.~\cite{Deng2025LDNF} investigated neighborhood fuzzy rough sets
for LDL feature selection, while Zhao et al.~\cite{Zhao2026IWNRS} developed
an informational weighted neighborhood rough set model and a conditional
discrimination index. These studies illustrate the relevance of granulation,
neighborhood construction and feature evaluation to label distribution data.

Granular-ball generation is also an important methodological issue.
The GBG++ method of Xie et al.~\cite{Xie2024GBG} addresses fast and stable
ball generation in classification. Its task setting differs from LDL,
but it provides relevant context for the design of splitting and termination
rules. Complementary directions include dynamic online LDL feature selection
using label importance and correlation~\cite{Chen2025OnlineLDL}, and
local model-agnostic explanations for LDL~\cite{Jia2025LIMEFLDL}.

Against this background, the present study investigates source-dependent
entropy weighting and directed dissimilarity together with attribute
acquisition costs. We distinguish the mathematical properties of the directed
dissimilarity from those of the adaptive granular-ball approximation model.
This distinction is necessary because changing the feature subset can change
the ball structure, so monotonicity of the dissimilarity does not alone
establish monotonicity of the resulting dependency measure.
